% !TEX program = xelatex
% Lernplatt - by Can Siebert
% Kurs 22 (Bo 143) - Tag 10 - Aufgabenheft
% Performance-Marketing im Vertrieb: Anzeigen, Conversion & E-Mail-Kampagnen
%
% Self-contained xelatex source. Kein biblatex, keine externen Bilder.

\documentclass[11pt,a4paper,oneside]{article}

% --- Encoding & Sprache --------------------------------------------------
\usepackage{polyglossia}
\setdefaultlanguage{german}

% --- Geometrie -----------------------------------------------------------
\usepackage[a4paper,margin=2.5cm]{geometry}

% --- Fonts ---------------------------------------------------------------
\usepackage{fontspec}
\defaultfontfeatures{Ligatures=TeX}

\IfFontExistsTF{Charter}{%
  \setmainfont{Charter}%
}{%
  \IfFontExistsTF{Noto Serif}{%
    \setmainfont{Noto Serif}%
  }{%
    \setmainfont{Liberation Serif}%
  }%
}

\IfFontExistsTF{Source Sans 3}{%
  \setsansfont{Source Sans 3}%
}{%
  \IfFontExistsTF{Source Sans Pro}{%
    \setsansfont{Source Sans Pro}%
  }{%
    \setsansfont{Liberation Sans}%
  }%
}

\newcommand{\caveatfontfallback}{\itshape}
\IfFontExistsTF{Caveat}{%
  \newfontfamily\caveatfont{Caveat}[Scale=1.15]%
}{%
  \newcommand\caveatfont{\caveatfontfallback}%
}

\setmonofont{Liberation Mono}

% --- Mikro-Typografie ----------------------------------------------------
\usepackage{microtype}
\usepackage{eurosym}
\linespread{1.15}

% --- Farben & Boxen ------------------------------------------------------
\usepackage{xcolor}
\definecolor{lpInk}{RGB}{20,28,38}
\definecolor{lpAccent}{RGB}{22,90,140}
\definecolor{lpAccentLight}{RGB}{210,228,240}
\definecolor{lpAmber}{RGB}{222,138,30}
\definecolor{lpAmberLight}{RGB}{255,243,217}
\definecolor{lpAmberBorder}{RGB}{196,118,18}
\definecolor{lpGray}{RGB}{96,104,114}
\definecolor{lpGrayLight}{RGB}{238,240,243}
\definecolor{lpGreen}{RGB}{34,120,72}
\definecolor{lpGreenLight}{RGB}{222,238,228}
\definecolor{lpGreenBorder}{RGB}{24,96,58}

% Niveau-Farben
\definecolor{lpL1}{RGB}{34,120,72}      % L1 = gruen (reproduktion)
\definecolor{lpL2}{RGB}{22,90,140}      % L2 = blau (anwendung)
\definecolor{lpL3}{RGB}{170,42,52}      % L3 = rot (management)

\usepackage[most]{tcolorbox}
\tcbuselibrary{breakable,skins}

% Info-Box (blau-weiss) fuer Hinweise/Aufgabenstellung-Rahmen
\newtcolorbox{infobox}[1][Hinweis]{%
  enhanced, breakable,
  colback=lpAccentLight,
  colframe=lpAccent,
  boxrule=0.6pt,
  arc=2pt,
  left=10pt,right=10pt,top=8pt,bottom=8pt,
  fonttitle=\sffamily\bfseries\color{white},
  coltitle=white,
  title={#1},
  attach boxed title to top left={xshift=8pt,yshift=-8pt},
  boxed title style={size=small, colback=lpAccent, colframe=lpAccent, boxrule=0.4pt, arc=1pt},
  top=14pt
}

% Antwort-Bereich: leerer Kasten mit Punktlinien
\newtcolorbox{antwortbox}[1][Ihre Antwort]{%
  enhanced, breakable,
  colback=white,
  colframe=lpGray,
  boxrule=0.4pt,
  arc=2pt,
  left=10pt,right=10pt,top=8pt,bottom=8pt,
  fonttitle=\sffamily\bfseries\color{white},
  coltitle=white,
  title={#1},
  attach boxed title to top left={xshift=8pt,yshift=-8pt},
  boxed title style={size=small, colback=lpGray, colframe=lpGray, boxrule=0.4pt, arc=1pt},
  top=14pt
}

% --- Listen --------------------------------------------------------------
\usepackage{enumitem}
\setlist{itemsep=2pt, topsep=4pt, parsep=0pt}
\setlist[itemize,1]{label={\textbullet},leftmargin=1.6em}
\setlist[itemize,2]{label={\textendash},leftmargin=1.4em}
\setlist[enumerate,1]{label=\arabic*., leftmargin=1.8em}

% --- Tabellen ------------------------------------------------------------
\usepackage{tabularx}
\usepackage{array}
\usepackage{booktabs}
\usepackage{longtable}

% --- Kopf-/Fusszeilen ----------------------------------------------------
\usepackage{fancyhdr}
\usepackage{lastpage}
\pagestyle{fancy}
\fancyhf{}
\renewcommand{\headrulewidth}{0.3pt}
\renewcommand{\footrulewidth}{0pt}
\fancyhead[L]{\footnotesize\sffamily\color{lpGray}Aufgabenheft \textperiodcentered{} Kurs 22 \textperiodcentered{} Tag 10}
\fancyhead[R]{\footnotesize\sffamily\color{lpGray}E-Mail \& Performance-Marketing}
\fancyfoot[L]{\footnotesize\sffamily\color{lpGray}\textcopyright{} Can Siebert}
\fancyfoot[R]{\footnotesize\sffamily\color{lpGray}\thepage{} / \pageref{LastPage}}

\fancypagestyle{plain}{%
  \fancyhf{}%
  \renewcommand{\headrulewidth}{0pt}%
  \fancyfoot[C]{\footnotesize\sffamily\color{lpGray}\thepage{} / \pageref{LastPage}}%
}

% --- Section-Styling -----------------------------------------------------
\usepackage[explicit]{titlesec}

\titleformat{\section}[block]
  {\sffamily\Large\bfseries\color{lpAccent}}
  {\thesection.}{0.6em}{#1}
  [{\vspace{2pt}\color{lpAccent}\titlerule[0.6pt]}]

\titleformat{\subsection}[block]
  {\sffamily\large\bfseries\color{lpInk}}
  {\thesubsection}{0.6em}{#1}

\titlespacing*{\section}{0pt}{14pt}{8pt}
\titlespacing*{\subsection}{0pt}{10pt}{4pt}

% --- Hyperlinks ----------------------------------------------------------
\usepackage[hidelinks,unicode]{hyperref}

% --- Helfer --------------------------------------------------------------
\newcommand{\akzent}[1]{{\caveatfont\color{lpAmber}#1}}
\newcommand{\fachbegriff}[1]{\textbf{#1}}

% Niveau-Badge
\newcommand{\niveaubadge}[2]{%
  \tcbox[on line, colback=#1, colframe=#1, coltext=white,
         boxrule=0pt, arc=2pt, left=5pt,right=5pt,top=1pt,bottom=1pt,
         nobeforeafter]{\sffamily\bfseries\footnotesize #2}%
}
\newcommand{\nivLi}{\niveaubadge{lpL1}{L1\,\textbullet\,Wissen}}
\newcommand{\nivLii}{\niveaubadge{lpL2}{L2\,\textbullet\,Anwendung}}
\newcommand{\nivLiii}{\niveaubadge{lpL3}{L3\,\textbullet\,Management}}

% Zeit-Hinweis
\newcommand{\zeit}[1]{%
  \tcbox[on line, colback=lpGrayLight, colframe=lpGrayLight, coltext=lpInk,
         boxrule=0pt, arc=2pt, left=5pt,right=5pt,top=1pt,bottom=1pt,
         nobeforeafter]{\sffamily\footnotesize\textbf{$\sim$\,#1}}%
}

% Aufgaben-Kopf
\newcommand{\aufgabenkopf}[4]{%
  \par\vspace{6pt}
  \noindent{\sffamily\Large\bfseries\color{lpAccent}Aufgabe #1 \textendash{} #2}\par
  \vspace{2pt}
  \noindent{\color{lpAccent}\rule{\linewidth}{0.6pt}}\par
  \vspace{4pt}
  \noindent #3 \hfill \zeit{#4}\par
  \vspace{6pt}
}

% YouTube-Suchquery
\newcommand{\ytquery}[1]{%
  \par\vspace{4pt}
  \noindent{\footnotesize\sffamily\color{lpGray}\textbf{Lernvideo zum Thema:}\ %
    \href{https://studyflix.de/search?query=#1}{\textcolor{lpAccent}{Studyflix-Treffer}}\ ·\ %
    \href{https://www.youtube.com/results?search\_query=#1}{\textcolor{lpAccent}{YouTube-Treffer}}\\%
    \textit{\textcolor{lpGray}{Stichworte: #1}}}\par
}

% Linierter Antwort-Bereich
\newcommand{\antwortlinien}[1]{%
  \par\vspace{6pt}
  \foreach \x in {1,...,#1}{%
    \noindent\makebox[\linewidth]{\dotfill}\\[6pt]
  }
}
\usepackage{pgffor}

% =========================================================================
\begin{document}

% --- COVER ---------------------------------------------------------------
\begin{titlepage}
  \pagestyle{empty}
  \vspace*{1.5cm}
  \noindent{\sffamily\footnotesize\color{lpGray}LERNPLATT \textperiodcentered{} BY CAN SIEBERT}\\[2pt]
  \noindent{\color{lpAccent}\rule{\linewidth}{1.2pt}}\\[4pt]
  \noindent{\sffamily\footnotesize\color{lpGray}AUFGABENHEFT \textperiodcentered{} MODUL 5 \textperiodcentered{} TAG 10 \textperiodcentered{} KURS 22 (Bo 143) \textperiodcentered{} 11.\,Juni 2026}

  \vspace{3cm}
  {\sffamily\fontsize{30}{34}\selectfont\bfseries\color{lpInk}Aufgabenheft\\Tag 10\par}

  \vspace{0.7cm}
  {\sffamily\large\color{lpGray}Performance-Marketing im Vertrieb:\\E-Mail-Kampagnen, Conversion\\\& ROI -- elf Aufgaben auf drei Niveaus.\par}

  \vspace{1.5cm}
  {\caveatfont\LARGE\color{lpAmber}11 Aufgaben \textperiodcentered{} L1 \textperiodcentered{} L2 \textperiodcentered{} L3}\par

  \vfill

  \noindent\begin{minipage}{0.55\linewidth}
    {\sffamily\footnotesize\color{lpGray}Niveau-Mix:}\\
    {\sffamily\small\color{lpInk}3\,\textsf{L1} \textperiodcentered{} 5\,\textsf{L2} \textperiodcentered{} 3\,\textsf{L3}}\\[10pt]
    {\sffamily\footnotesize\color{lpGray}Bearbeitung:}\\
    {\sffamily\small\color{lpInk}Selbstbearbeitung im Lernpfad}
  \end{minipage}\hfill
  \begin{minipage}{0.4\linewidth}
    \raggedleft
    {\sffamily\footnotesize\color{lpGray}Dozent:}\\
    {\sffamily\small\color{lpInk}Can Siebert}\\[10pt]
    {\sffamily\footnotesize\color{lpGray}Begleitend:}\\
    {\sffamily\small\color{lpInk}Lehrskript \& Lösungsheft}
  \end{minipage}

  \vspace{0.5cm}
  \noindent{\color{lpAccent}\rule{\linewidth}{0.6pt}}
\end{titlepage}

% --- ANLEITUNG -----------------------------------------------------------
\section*{So arbeiten Sie mit diesem Heft}
\thispagestyle{plain}

\noindent Dieses Aufgabenheft enthält elf Aufgaben zu Tag 10 \textendash{} \emph{Performance-Marketing im Vertrieb: E-Mail-Kampagnen, Conversion \& ROI}. Die Aufgaben sind in drei Niveaus gegliedert:

\begin{itemize}
  \item \nivLi{} \textendash{} \fachbegriff{Wissen.} Kurze Reproduktion und Einordnung. 5\,--\,10 Minuten pro Aufgabe.
  \item \nivLii{} \textendash{} \fachbegriff{Anwendung.} Transfer auf konkrete Vertriebssituationen. 15\,--\,30 Minuten.
  \item \nivLiii{} \textendash{} \fachbegriff{Management.} Entscheidungsarchitektur und Verantwortung. 30\,--\,45 Minuten.
\end{itemize}

\noindent Jede Aufgabe enthält eine Niveau-Markierung, einen Zeit-Richtwert, den Aufgabentext (oft mit Kontext oder Fallbeschreibung), einen Antwort-Freiraum für Ihre Lösung sowie eine \emph{YouTube-Suchquery} für den begleitenden Lernpfad.

\begin{infobox}[Hinweis]
Die Musterlösungen finden Sie im separaten \emph{Lösungsheft Tag 10}. Versuchen Sie jede Aufgabe zuerst eigenständig \textendash{} der Mehrwert entsteht in der eigenen Argumentation, nicht im Abschreiben.
\end{infobox}

\clearpage

% =========================================================================
% L1 AUFGABEN
% =========================================================================
\section*{Niveau L1 \textendash{} Wissen \& Reproduktion}
\addcontentsline{toc}{section}{L1 -- Wissen}

% --- AUFGABE 1 -----------------------------------------------------------
% LERNPLATT_HILFSVIDEOS
\section*{Hilfsvideos zu diesem Tag}
\addcontentsline{toc}{section}{Hilfsvideos zu diesem Tag}
\thispagestyle{plain}

\noindent Die folgenden YouTube-Erkl\"arvideos vermitteln die Inhalte, die Sie zur Bearbeitung der Aufgaben ben\"otigen. Klicken Sie auf den Titel, um das Video direkt zu \"offnen; die vollst\"andige URL steht in Klammern.

\begin{description}[style=nextline,leftmargin=18pt,labelindent=0pt,itemsep=8pt]
  \item[1.~E-Mail-Marketing als systematischer Vertriebshebel] \href{https://youtu.be/GlKFLDx0Fjw}{\textcolor{lpAccent}{https://youtu.be/GlKFLDx0Fjw}}\\[2pt]
  {\footnotesize\color{lpGray}(URL: \nolinkurl{https://youtu.be/GlKFLDx0Fjw})}
  \item[2.~Lead Nurturing mit Marketing Automation] \href{https://youtu.be/sjAkl8pwLgg}{\textcolor{lpAccent}{https://youtu.be/sjAkl8pwLgg}}\\[2pt]
  {\footnotesize\color{lpGray}(URL: \nolinkurl{https://youtu.be/sjAkl8pwLgg})}
  \item[3.~Newsletter erfolgreich starten — Aufbau, Segmentierung, Trigger] \href{https://youtu.be/NPGwFJqX0kc}{\textcolor{lpAccent}{https://youtu.be/NPGwFJqX0kc}}\\[2pt]
  {\footnotesize\color{lpGray}(URL: \nolinkurl{https://youtu.be/NPGwFJqX0kc})}
  \item[4.~Kampagnen und Automatisierung in der Praxis] \href{https://youtu.be/RF2m7mwJmxc}{\textcolor{lpAccent}{https://youtu.be/RF2m7mwJmxc}}\\[2pt]
  {\footnotesize\color{lpGray}(URL: \nolinkurl{https://youtu.be/RF2m7mwJmxc})}
\end{description}
\vspace{0.6em}
\noindent{\footnotesize\color{lpGray}\textit{Tipp:} \"Offnen Sie das Video, lassen Sie es im Hintergrund laufen und arbeiten Sie parallel an der Aufgabe.}

\clearpage

\aufgabenkopf{1}{Newsletter \textendash{} Kampagne \textendash{} Nurturing \textendash{} Vertriebsstrecke}{\nivLi}{8\,min}

\noindent Definieren Sie in jeweils \emph{einem} Satz die vier Begriffe, die im professionellen E-Mail-Vertrieb häufig vermischt werden:

\begin{enumerate}
  \item Newsletter
  \item Kampagne
  \item Lead-Nurturing
  \item Automatisierte Vertriebsstrecke
\end{enumerate}

\noindent Ergänzen Sie einen fünften Satz: Worin liegt der zentrale Unterschied zwischen einem \emph{Newsletter} und einer \emph{Nurturing-Strecke} aus Vertriebssicht?

\ytquery{Newsletter Lead Nurturing E-Mail Marketing Unterschied B2B Vertrieb}

\antwortlinien{10}

\clearpage

% --- AUFGABE 2 -----------------------------------------------------------
\aufgabenkopf{2}{E-Mail-Kennzahlen zuordnen}{\nivLi}{8\,min}

\noindent Ordnen Sie jede der folgenden Kennzahlen einer der drei Kategorien zu \textendash{} und schreiben Sie in einem Halbsatz, warum die Zuordnung sinnvoll ist:

\medskip
\noindent\textbf{Kategorien:} (A) Aktivitätskennzahl \textperiodcentered{} (B) Qualitätskennzahl \textperiodcentered{} (C) Wirtschaftliche Steuerungskennzahl.

\medskip
\noindent\textbf{Kennzahlen:}

\begin{enumerate}
  \item Öffnungsrate
  \item Klickrate
  \item Conversion Rate zu Beratungstermin
  \item Abmelderate
  \item Anteil vertriebsqualifizierter Leads
  \item Kosten pro qualifiziertem Lead
  \item Deckungsbeitrag aus E-Mail-Kampagne
  \item Bounce Rate
\end{enumerate}

\ytquery{E-Mail Marketing KPI Öffnungsrate Conversion Pipeline Beitrag ROI}

\antwortlinien{12}

\clearpage

% --- AUFGABE 3 -----------------------------------------------------------
\aufgabenkopf{3}{ROI, CAC und Conversion Rate auf den Punkt}{\nivLi}{8\,min}

\noindent Drei wirtschaftliche Kerngrößen werden im Performance-Marketing regelmäßig verwendet \textendash{} und ebenso regelmäßig verwechselt.

\begin{enumerate}
  \item Erklären Sie in \emph{einem} Satz, was der \emph{Return on Investment (ROI)} im Vertrieb misst.
  \item Erklären Sie in \emph{einem} Satz, was \emph{Customer Acquisition Cost (CAC)} bedeutet.
  \item Erklären Sie in \emph{einem} Satz, was eine \emph{Conversion Rate} zum Beratungstermin aussagt \textendash{} und worin sie sich von einer Klickrate unterscheidet.
\end{enumerate}

\noindent Ergänzen Sie einen vierten Satz: Welche dieser drei Kennzahlen ist am ehesten eine \emph{Managementkennzahl}, welche eher eine \emph{operative Steuerungskennzahl}?

\ytquery{ROI CAC Conversion Rate Definition Vertrieb Performance Marketing}

\antwortlinien{10}

\clearpage

% =========================================================================
% L2 AUFGABEN
% =========================================================================
\section*{Niveau L2 \textendash{} Anwendung \& Transfer}
\addcontentsline{toc}{section}{L2 -- Anwendung}

% --- AUFGABE 4 -----------------------------------------------------------
\aufgabenkopf{4}{CloudDesk Solutions \textendash{} Nurturing-Strecke ohne Vorwissen}{\nivLii}{25\,min}

\begin{infobox}[Fallbeschreibung]
\textbf{Unternehmen:} ``CloudDesk Solutions'' bietet eine Softwarelösung für digitales Dokumentenmanagement im Mittelstand an. Auf der Website können Interessenten ein Whitepaper herunterladen: ``5 Wege zu effizienteren Dokumentenprozessen''. Viele Personen laden das Whitepaper herunter, aber nur wenige buchen anschließend einen Beratungstermin.

\smallskip
\textbf{Gegeben:}
\begin{itemize}
  \item Zielgruppe: Geschäftsführer, IT-Leiter, kaufmännische Leiter
  \item Produkt ist erklärungsbedürftig
  \item Verkaufszyklus: 2 bis 4 Monate
  \item Vertrieb möchte nur qualifizierte Leads erhalten
  \item Marketing hat bisher nur eine automatische Danke-Mail eingerichtet
  \item Ziel: mehr Beratungstermine aus bestehenden Whitepaper-Leads gewinnen
\end{itemize}
\end{infobox}

\noindent\textbf{Arbeitsauftrag:}

\begin{enumerate}
  \item Beschreiben Sie in 2\,--\,3 Sätzen, welches Problem im aktuellen Prozess besteht.
  \item Benennen Sie \textbf{drei mögliche Gründe}, warum Whitepaper-Leads noch keinen Beratungstermin buchen.
  \item Entwickeln Sie eine einfache E-Mail-Strecke mit \textbf{mindestens vier E-Mails} nach dem Download.
  \item Legen Sie für jede E-Mail \emph{Ziel}, \emph{Inhalt} und \emph{Call-to-Action} fest.
  \item Entscheiden Sie, \emph{wann} ein Lead an den Vertrieb übergeben werden sollte.
\end{enumerate}

\ytquery{Lead Nurturing E-Mail Strecke Whitepaper B2B Beratungstermin}

\antwortlinien{20}

\clearpage

% --- AUFGABE 5 -----------------------------------------------------------
\aufgabenkopf{5}{E-Mail-Performance analysieren und optimieren}{\nivLii}{25\,min}

\begin{infobox}[Kennzahlen der bisherigen Kampagne]
\textbf{CloudDesk Solutions \textendash{} aktuelle Whitepaper-Kampagne:}
\begin{itemize}
  \item 1.200 Whitepaper-Downloads
  \item Danke-Mail Öffnungsrate: 62\,\%
  \item Klickrate auf Beratungslink: 4\,\%
  \item Beratungstermine: 18
  \item Davon qualifizierte Termine: 7
  \item Abmelderate: 3,8\,\%
  \item Vertrieb meldet: ``Viele Leads sind noch nicht entscheidungsbereit.''
\end{itemize}

\smallskip
\textbf{Rahmenbedingungen:}
\begin{itemize}
  \item Budget für Optimierung: 18.000\,\euro
  \item Zeitraum: 3 Monate
  \item Ziel: 40 qualifizierte Beratungstermine
  \item Vertrieb kann maximal 15 neue Termine pro Woche bearbeiten
\end{itemize}
\end{infobox}

\noindent\textbf{Arbeitsauftrag:}

\begin{enumerate}
  \item Bewerten Sie die aktuelle Kampagne anhand der Kennzahlen \textendash{} was ist schon gut, was bleibt zurück?
  \item Identifizieren Sie \textbf{mindestens vier Schwachstellen} im aktuellen Nurturing-Prozess.
  \item Entwickeln Sie \textbf{drei Optimierungsmaßnahmen} für Inhalte, Timing oder Segmentierung.
  \item Legen Sie \textbf{fünf KPIs} fest, mit denen die neue E-Mail-Strecke gesteuert werden soll.
  \item Entscheiden Sie, ob das Ziel von 40 qualifizierten Beratungsterminen realistisch ist. Begründen Sie Ihre Einschätzung.
\end{enumerate}

\ytquery{E-Mail Kampagne Performance Analyse Optimierung Conversion B2B Lead}

\antwortlinien{18}

\clearpage

% --- AUFGABE 6 -----------------------------------------------------------
\aufgabenkopf{6}{Segmentierung einer Kontaktbasis}{\nivLii}{20\,min}

\begin{infobox}[Ausgangslage]
\textbf{Software-Anbieter mit 8.000 Kontakten:} Webinar-Teilnehmende, Whitepaper-Downloads, LinkedIn-Leads, Messekontakte, inaktive Altkontakte. Bisher erhalten \emph{alle} Kontakte denselben monatlichen Newsletter. Die Geschäftsführung fragt: ``Wie würden wir die Liste sinnvoll segmentieren \textendash{} und welche Daten brauchen wir dafür?''
\end{infobox}

\noindent\textbf{Arbeitsauftrag:}

\noindent Entwickeln Sie eine Segmentierungslogik entlang \textbf{vier Dimensionen}. Für jede Dimension:

\begin{itemize}
  \item Konkrete Segmente (mindestens drei pro Dimension).
  \item Welche Information aus dem CRM oder dem Tracking ist nötig?
  \item Welcher Inhaltstyp passt zu welchem Segment?
\end{itemize}

\noindent Vorgeschlagene Dimensionen: \emph{Rolle / Buying-Center}, \emph{Branche}, \emph{Interesse / Themenpräferenz}, \emph{Reifegrad / Kaufphase}.

\smallskip
\noindent Formulieren Sie zum Schluss \emph{eine} Hypothese: Welches Segment hat den höchsten Hebel auf qualifizierte Vertriebsleads in den nächsten 90 Tagen?

\ytquery{B2B E-Mail Segmentierung Rolle Branche Reifegrad CRM Nurturing}

\antwortlinien{16}

\clearpage

% --- AUFGABE 7 -----------------------------------------------------------
\aufgabenkopf{7}{EcoFleet Systems \textendash{} Nurturing- und KPI-Architektur}{\nivLii}{30\,min}

\begin{infobox}[Fallstudie: EcoFleet Systems]
``EcoFleet Systems'' bietet digitale Lösungen zur Fuhrparkoptimierung für mittelständische Unternehmen an. Die Software senkt Fahrzeugkosten, reduziert CO\textsubscript{2}-Emissionen und macht Wartungsprozesse transparenter. Bisher gewinnt das Unternehmen Leads über Webinare, Whitepaper, LinkedIn-Kampagnen und Messekontakte. Viele Kontakte zeigen Interesse, nur ein kleiner Teil entwickelt sich zur qualifizierten Verkaufschance.

\smallskip
\textbf{Rahmen:} 12\,Mio.\,\euro{} Jahresumsatz \textperiodcentered{} 7 Vertriebler \textperiodcentered{} 2 im Marketing \textperiodcentered{} 8.000 Kontakte (davon 2.500 Webinar-/Whitepaper-Leads) \textperiodcentered{} Auftragswert 38.000\,\euro/Jahr \textperiodcentered{} Verkaufszyklus 3\,--\,6 Monate \textperiodcentered{} Budget 95.000\,\euro{} für 6\,Monate \textperiodcentered{} Ziel: 300 qualifizierte Leads, 90 Demo-Termine, 25 Angebote, 8 Abschlüsse.

\smallskip
\textbf{Probleme:} alle Kontakte erhalten gleiche E-Mails, keine Segmentierung, Vertrieb erhält Leads zu früh, Bewertung läuft über Öffnungs-/Klickraten, CRM-Daten uneinheitlich, kein gemeinsames Lead-Verständnis.
\end{infobox}

\noindent\textbf{Arbeitsauftrag:}

\begin{enumerate}
  \item Analysieren Sie die \textbf{drei wichtigsten Schwächen} des aktuellen E-Mail- und Leadprozesses.
  \item Entwickeln Sie eine \textbf{Nurturing-Strecke mit mindestens sechs Kontaktpunkten} entlang der Customer Journey.
  \item Definieren Sie, wann ein Lead \emph{marketingqualifiziert (MQL)} und wann \emph{vertriebsqualifiziert (SQL)} ist.
  \item Bauen Sie eine \textbf{KPI-Architektur} von Zustellung bis Umsatz auf (mindestens sechs Kennzahlen).
\end{enumerate}

\ytquery{EcoFleet B2B Lead Nurturing KPI Architektur MQL SQL Marketing Automation}

\antwortlinien{22}

\clearpage

% --- AUFGABE 8 -----------------------------------------------------------
\aufgabenkopf{8}{Budgetverteilung 95.000\,\euro{} \textendash{} mit Begründung}{\nivLii}{20\,min}

\noindent\textbf{Arbeitsauftrag:}

\noindent Bleiben Sie im EcoFleet-Szenario. Sie haben 95.000\,\euro{} für sechs Monate. Verteilen Sie das Budget auf die folgenden Posten \textendash{} und begründen Sie jede Zuteilung in einem Satz:

\medskip
\noindent\begin{tabularx}{\linewidth}{@{}lXl@{}}
\toprule
\textbf{Posten} & \textbf{Begründung in einem Satz} & \textbf{Betrag (\euro)} \\
\midrule
Contententwicklung (Cases, Checklisten, Fachbeiträge) & \rule{8em}{0.3pt} & \rule{3em}{0.3pt} \\
\midrule
Marketing-Automation und technische Einrichtung & \rule{8em}{0.3pt} & \rule{3em}{0.3pt} \\
\midrule
CRM-Datenbereinigung und Segmentierung & \rule{8em}{0.3pt} & \rule{3em}{0.3pt} \\
\midrule
Landing Pages, Demo-Formulare, Conversion-Optimierung & \rule{8em}{0.3pt} & \rule{3em}{0.3pt} \\
\midrule
A/B-Tests (Betreff, Inhalt, CTA, Timing) & \rule{8em}{0.3pt} & \rule{3em}{0.3pt} \\
\midrule
Reporting, Dashboards, ROI-Auswertung & \rule{8em}{0.3pt} & \rule{3em}{0.3pt} \\
\midrule
\textbf{Summe} & & \textbf{95.000} \\
\bottomrule
\end{tabularx}

\medskip
\noindent Schließen Sie mit einer \textbf{Empfehlung in drei Sätzen}: Welche zwei Posten sind aus Ihrer Sicht die wichtigsten \textendash{} und welcher Posten könnte am ehesten reduziert werden, falls das Budget gekürzt würde?

\ytquery{Marketing Budget B2B Nurturing CRM Automation Content Verteilung}

\antwortlinien{14}

\clearpage

% =========================================================================
% L3 AUFGABEN
% =========================================================================
\section*{Niveau L3 \textendash{} Management \& Entscheidungsarchitektur}
\addcontentsline{toc}{section}{L3 -- Management}

% --- AUFGABE 9 -----------------------------------------------------------
\aufgabenkopf{9}{EcoFleet \textendash{} Entscheidung unter Unsicherheit}{\nivLiii}{40\,min}

\begin{infobox}[Rolle und Spannungsfeld]
\textbf{Ihre Rolle:} Chief Revenue Officer von EcoFleet Systems (siehe Aufgabe 7).

\smallskip
\textbf{Spannungsfeld:} Marketing will mit dem 95.000-\euro{}-Budget primär \emph{mehr Leads} generieren (LinkedIn-Ads, Webinare, Whitepaper-Kampagnen). Vertrieb sagt: ``Wir haben schon zu viele unreife Leads \textendash{} schickt uns lieber weniger, aber bessere.'' Geschäftsführung will harte Zahlen: \emph{8 Abschlüsse, 300 qualifizierte Leads}.

\smallskip
\textbf{Sie müssen entscheiden:} Investieren Sie zuerst in \emph{Lead-Quantität} (mehr Kontakte aus neuen Kanälen) oder in \emph{Lead-Qualität} (Segmentierung, Datenqualität, Scoring, Nurturing-Inhalte für die 2.500 bestehenden Kontakte)?
\end{infobox}

\noindent\textbf{Arbeitsauftrag:}

\begin{enumerate}
  \item Analysieren Sie das Spannungsfeld in 3\,--\,5 Punkten. Welche Annahmen liegen hinter ``mehr Leads'' und welche hinter ``bessere Leads''?
  \item Entwickeln Sie \textbf{zwei strategische Optionen} (A: Quantität zuerst, B: Qualität zuerst) und beschreiben Sie deren Investitionslogik.
  \item Bewerten Sie beide Optionen anhand von: \emph{Wahrscheinlichkeit der Zielerreichung} \textperiodcentered{} \emph{Risiko für Vertriebskapazität} \textperiodcentered{} \emph{ROI-Sichtbarkeit nach 6 Monaten} \textperiodcentered{} \emph{Datenkontrolle} \textperiodcentered{} \emph{Skalierbarkeit ab Monat 7}.
  \item Treffen Sie eine begründete \textbf{Entscheidung} und legen Sie die Budget-Aufteilung zwischen Quantität und Qualität fest.
  \item Beschreiben Sie, welche \textbf{Pilot-Logik} Sie nutzen, um Ihre Entscheidung in den ersten 8 Wochen zu validieren \textendash{} und an welcher Kennzahl Sie sie ggf.\ korrigieren würden.
  \item Formulieren Sie \emph{eine ausdrückliche Entscheidung unter Unsicherheit}: An welcher Stelle entscheiden Sie ohne ausreichende Datenlage \textendash{} und warum trotzdem so?
\end{enumerate}

\ytquery{Lead Qualität Quantität B2B Vertrieb CRO Entscheidung Pipeline Pilot}

\antwortlinien{32}

\clearpage

% --- AUFGABE 10 ----------------------------------------------------------
\aufgabenkopf{10}{Zielkonflikt: Automatisierung vs.\ Personalisierung}{\nivLiii}{30\,min}

\noindent\textbf{Diskussionsaufgabe.}

\noindent Im professionellen E-Mail-Vertrieb stehen \emph{Automatisierung} (Skalierung, Effizienz, Wirtschaftlichkeit) und \emph{Personalisierung} (Relevanz, Kundenerlebnis, Conversion) in einem dauerhaften Spannungsverhältnis. Datenschutz, Vertriebslast und Kundenwahrnehmung kommen hinzu.

\medskip
\noindent\textbf{Arbeitsauftrag:}

\begin{enumerate}
  \item Beschreiben Sie in eigenen Worten, ab \emph{wann} Automatisierung in der Kundenwahrnehmung kippt: vom hilfreichen Service zur ``Massen-Mail-Belästigung'' (3\,--\,4 Sätze). Was sind die kritischen Schwellen?
  \item Wählen Sie ein \textbf{konkretes B2B-Szenario} (z.\,B.\ SaaS-Anbieter mit 50.000 Kontakten, Industriedienstleister mit Bestandskundenliste, mittelständischer Maschinenbauer mit Webinar-Leads). Erläutern Sie:
  \begin{itemize}
    \item Welche \emph{drei Automatisierungsmöglichkeiten} wären hier sinnvoll \textendash{} und welche eine Grenze überschreiten würde?
    \item Welche \emph{drei Personalisierungs-Hebel} sind ohne überzogenen Aufwand realistisch?
    \item Wo entstehen Konflikte mit Datenschutz, Einwilligung, Vertriebskapazität oder Markenwahrnehmung?
  \end{itemize}
  \item Formulieren Sie eine \textbf{Faustregel} (ein bis zwei Sätze): Wann ist Automatisierung wirtschaftlich begründet \textendash{} und wann beginnt sie, mehr Vertrauen zu zerstören als Umsatz zu erzeugen?
\end{enumerate}

\ytquery{Marketing Automation Personalisierung B2B Grenzen Datenschutz DSGVO}

\antwortlinien{22}

\clearpage

% --- AUFGABE 11 ----------------------------------------------------------
\aufgabenkopf{11}{Transferprojekt \textendash{} E-Mail-, Nurturing- \& ROI-Architektur}{\nivLiii}{45\,min}

\begin{infobox}[Rolle und Ausgangslage]
\textbf{Ihre Rolle:} Chief Revenue Officer \emph{oder} Head of Marketing Automation \emph{oder} Leiter:in Digital Sales eines mittelständischen B2B-Unternehmens.

\smallskip
\textbf{Ausgangslage:} Wachsende Kontaktdatenbank aus Webinaren, Whitepapers, Plattformkontakten, Messen, LinkedIn-Kampagnen und Bestandskunden. Kontakte werden unregelmäßig mit Newslettern bespielt. Keine Nurturing-Architektur, keine verbindliche Leadqualifizierung, keine belastbare ROI-Messung. Vertrieb beklagt unreife Leads, Marketing beklagt fehlendes Feedback, Geschäftsführung fordert wirtschaftliche Nachweise.

\smallskip
\textbf{Rahmenbedingungen:}
\begin{itemize}
  \item Budget begrenzt
  \item Datenqualität uneinheitlich
  \item Vertriebskapazität knapp
  \item Kaufentscheidungen dauern mehrere Monate
  \item Datenschutz und Einwilligungen zwingend zu beachten
  \item Geschäftsführung erwartet belastbare ROI-Aussagen
  \item Marketing und Vertrieb nutzen unterschiedliche Erfolgskennzahlen
\end{itemize}
\end{infobox}

\noindent\textbf{Arbeitsauftrag:} Entwickeln Sie eine strukturierte \emph{Entscheidungsarchitektur} \textendash{} keine reine Maßnahmenliste \textendash{} für E-Mail-Marketing, Lead-Nurturing und ROI-Optimierung.

\begin{enumerate}
  \item Zielbild für die E-Mail- und Nurturing-Architektur (3\,--\,5 Sätze).
  \item Segmentierung der Kontaktbasis nach Zielgruppe, Interesse, Kaufphase und Verhalten.
  \item Nurturing-Logik entlang der Customer Journey (Erstkontakt bis Wiederansprache).
  \item Inhaltsarchitektur: E-Mail-Typen, Themen, Trigger, Call-to-Actions.
  \item Lead-Scoring- und Übergabelogik an den Vertrieb (MQL/SQL-Definition).
  \item Daten- und Datenschutzlogik (Einwilligung, CRM, Aktualität, Qualität).
  \item KPI-Architektur von Zustellung bis Umsatz und Deckungsbeitrag.
  \item ROI-Logik zur Bewertung von Kampagnen, Segmenten und Vertriebsbeitrag.
  \item Governance zwischen Marketing, Vertrieb, CRM, Datenschutz, Controlling, Geschäftsführung.
  \item Drei Managemententscheidungen unter Unsicherheit.
\end{enumerate}

\noindent\textbf{Zusatzanforderung:} Treffen Sie mindestens \emph{eine bewusste Entscheidung gegen mehr Kampagnenvolumen}, wenn dadurch Leadqualität, Kundenerlebnis, Datenschutz, Vertriebskapazität oder Wirtschaftlichkeit gefährdet würden \textendash{} und begründen Sie diese aus Senior-Level-Perspektive.

\ytquery{Marketing Automation Architektur Lead Nurturing ROI Governance CRO}

\antwortlinien{36}

\clearpage

% --- Abschluss -----------------------------------------------------------
\section*{Abschluss}
\thispagestyle{plain}

\noindent Sie haben alle elf Aufgaben bearbeitet. Vergleichen Sie nun Ihre Antworten mit dem \emph{Lösungsheft Tag 10}.

\medskip
\begin{infobox}[Was Sie jetzt können sollten]
\begin{itemize}
  \item Newsletter, Kampagne, Nurturing und automatisierte Vertriebsstrecke sauber abgrenzen.
  \item E-Mail-Kennzahlen nach Aktivität, Qualität und wirtschaftlicher Wirkung sortieren.
  \item Eine Nurturing-Strecke entlang der Customer Journey entwerfen.
  \item Lead-Qualifizierung (MQL/SQL) verbindlich definieren.
  \item Eine KPI- und ROI-Architektur aufbauen, die Marketing- und Vertriebslogik verbindet.
  \item Entscheidungen unter Unsicherheit zwischen Lead-Quantität und Lead-Qualität verantworten.
  \item Eine Architektur formulieren, in der Automatisierung, Personalisierung, Datenschutz und Wirtschaftlichkeit balanciert sind.
\end{itemize}
\end{infobox}

\vspace{1cm}
\noindent{\color{lpAccent}\rule{\linewidth}{0.6pt}}\\[4pt]
{\footnotesize\sffamily\color{lpGray}Lernplatt \textperiodcentered{} by Can Siebert \textperiodcentered{} Kurs 22 (Bo 143) \textperiodcentered{} Tag 10 \textperiodcentered{} Aufgabenheft \textperiodcentered{} \textcopyright{} 2026}

\end{document}
